% fig-d6-causal-spine.tex --- D6, Paper III's causal spine.
%
% Defines \ClebschCausalSpine[<variant>], where <variant> is "current" or
% "final".  The argument is optional and defaults to "current", so a bare
% \ClebschCausalSpine keeps drawing what it always drew.  Any unrecognised
% argument is treated as "current".
%
% Left to right in both variants: the chain of the Paper III abstract, from
% the incidence cover to the operator's shadows.  Only the return arrow
% underneath differs.
%
%   current --- the recognition statement the manuscript actually proves
%     (thm:aligned-faithfulness): marked determinant-(-3) four-subsets
%     recover the conference signing up to switching and global negation.
%
%   final --- the triangle/Pfaffian recognition of
%     notes/2026-08-03-c862-paper-iii-spectral-descent-recognition-theorem-packet.md
%     (Theorem III there, proved earlier in C809).  Two hypotheses do two
%     different jobs and must stay separate in the label: all off-diagonal
%     entries nonzero is what forces A^2 = lambda I on the whole nonzero-edge
%     torus, and the scalar sign locus is what pins lambda = 5 and produces
%     the unique order-six conference class with its two outer orientations.
%
% The return arrow of the "final" variant is drawn with csfforth, a DOUBLE
% rule, and tagged with a csfbadge reading "proved; awaiting promotion".
% This is a third claim status and must not be drawn like either of the two
% D1 caveats: the dashed csfopen edge there marks a bridge no theorem
% supports, and the Paper V node marks a paper that does not exist.  This
% theorem is proved; only its appearance in the manuscript is pending.  The
% Lean note under it is deliberately narrower than the human statement,
% because the formalisation is: it is stated over an integral domain, and
% its orientation split rests on one compiled finite check rather than on a
% kernel-checked parity argument
% (notes/2026-08-02-c815-four-shadow-lean-formalization.md).

\providecommand{\ClebschCausalSpine}[1][current]{%
\csfSetSpineReturn{#1}%
\begin{tikzpicture}[
  spine/.style={csfnode,font=\footnotesize,text width=19mm,
    minimum height=17mm,anchor=north},
  every node/.style={outer sep=0pt},
]

\node[spine] (c1) at (-6.0,0)
  {incidence cover\\[1pt] degree two over \(\mathbf P(H)\)};
\node[spine] (c2) at (-3.6,0)
  {residue-field pinching\\[1pt] nonsplit two-branch};
\node[spine] (c3) at (-1.2,0)
  {square class \([5]\)\\[1pt] \(\sqrt{5J_{0}}\)};
\node[spine] (c4) at (1.2,0)
  {marked bridge datum\\[1pt] fixes the sign};
\node[spine,line width=1pt,fill=black!8] (c5) at (3.6,0)
  {conference operator\\[1pt] \(C\), \(C^{2}=5I\)};
\node[spine] (c6) at (6.0,0)
  {shadows\\[1pt] \(\operatorname{Pf}[D_x,C]\), \(\det\), norm};

\foreach \a/\b in {c1/c2,c2/c3,c3/c4,c4/c5,c5/c6}
  {\draw[csfarr] (\a.east) -- (\b.west);}

\ifnum\csfSpineFinal=1

\draw[csfforth] (c6.south) -- (6.0,-2.35) -- (3.6,-2.35) -- (c5.south);
\node[csfbadge,anchor=east] at (3.25,-2.35) {proved; awaiting promotion};

\node[anchor=north,align=flush center,font=\scriptsize,text width=120mm]
  (spinestmt) at (0,-2.95)
  {recognition, stronger form: for symmetric \(A\) with zero diagonal and
   \emph{every} off-diagonal entry nonzero, \(\operatorname{Pf}[D_x,A]=\mu
   T_A(x)\) with \(\mu\ne0\) forces \(A^{2}=\lambda I\); restricted to the
   scalar sign locus \(A=sB\), \(b_{ij}=\pm1\), it recovers the unique
   order-six conference class together with its two outer orientations
   \(\mu=\pm4\)};

\node[align=flush center,font=\scriptsize,text width=110mm,
  below=3.5mm of spinestmt] (spinedeg)
  {The order is a conclusion, not a hypothesis:
   \(\operatorname{Pf}[D_x,A]\) has degree \(n/2\) for an operator of order
   \(n\), while the triangle cubic \(T_A\) has degree three, so
   unhomogenized proportionality is possible only at \(n=6\).};

\node[align=flush center,font=\tiny,text width=110mm,
  below=3.5mm of spinedeg]
  {Proved in characteristic \(\ne2\), and formalized over an integral
   domain: the quadratic converse and the pentagon gauge are kernel-checked,
   while the orientation split rests on one compiled finite check.  It is
   not yet a statement of the manuscript.};

\else

\draw[csfarr] (c6.south) -- (6.0,-2.35) -- (3.6,-2.35) -- (c5.south);
\node[anchor=north,align=center,font=\scriptsize,text width=90mm]
  at (2.4,-2.55)
  {recognition: the marked determinant-\((-3)\) four-subsets recover the
   signing up to switching and global negation};

\fi

\end{tikzpicture}%
}
